\chapter{2012年湖北卷（理）}

\section{单选题}

\begin{problem}
方程 \(x^{2} + 6x + 13 = 0\) 的一个根是（\quad）

\choices
    {\(-3 + 2i\)}
    {\(3 + 2i\)}
    {\(-2 + 3i\)}
    {\(2 + 3i\)}
\end{problem}

\begin{problem}
命题“\(\exists x_0 \in \complement_B Q, x_0^2 \in Q\)”的否定是（\quad）

\choices
    {\(\exists x_0 \notin \complement_B Q, x_0^2 \in Q\)}
    {\(\exists x_0 \in \complement_B Q, x_0^2 \notin Q\)}
    {\(\forall x \notin \complement_B Q, x^3 \in Q\)}
    {\(\forall x \in \complement_B Q, x^3 \notin Q\)}
\end{problem}

\begin{problem}
已知二次函数 \( y = f(x) \) 的图象如图所示，则它与 \( x \) 轴所围成的图形的面积为（\quad）

\choices
    {\( \frac{2\pi}{5} \)}
    {\( \frac{4}{3} \)}
    {\( \frac{3}{2} \)}
    {\( \frac{\pi}{2} \)}
\end{problem}

\begin{problem}
已知某几何体的三视图如图所示，则该几何体的体积为（\quad）

\bitmapfigure[max width=0.74\linewidth,max totalheight=5.8cm]{img/2012/hubei_science/q4_fig1.png}

\choices
    {\(\frac{8\pi}{3}\)}
    {\(3\pi\)}
    {\(\frac{10\pi}{3}\)}
    {\(6\pi\)}
\end{problem}

\begin{problem}
设 \(a \in \Z\)，且 \(0 \leqslant a < 13\)，若 \(5^{1.2012} + a\) 能被 13 整除，则 \(a =\)（\quad）

\choices
    {0}
    {1}
    {11}
    {12}
\end{problem}

\begin{problem}
设 \(a, b, c, x, y, z\) 是正数，且 \(a^2 + b^2 + c^2 = 10, x^2 + y^2 + z^2 = 40\)，\(\frac{a + b + c}{x + y + z} = 20\)，则 \(\frac{a + b + c}{x + y + z} =\)（\quad）

\choices
    {\( \frac{1}{4} \)}
    {\( \frac{1}{3} \)}
    {\( \frac{1}{2} \)}
    {\( \frac{3}{4} \)}
\end{problem}

\begin{problem}
定义在 \((- \infty, 0) \cup (0, +\infty)\) 上的函数 \(f(x)\)，如果对于任意给定的等比数列 \(\{a_n\}\)，\(\{f(a_n)\}\) 仍是等比数列，则称 \(f(x)\) 为“保等比数列函数”. 现有定义在 \((- \infty, 0) \cup (0, +\infty)\) 上的如下函数: \circled{1} \( f(x)=x^{2} \) ; \circled{2} \( f(x)=2^{x} \) ; \circled{3} \( f(x)=\sqrt{|x|} \) ; \circled{4} \( f(x)=\ln|x| \) . 则其中是“保等比数列函数”的 \(f(x)\) 的序号为（\quad）

\choices
    {\circled{1}\circled{2}}
    {\circled{3}\circled{4}}
    {\circled{3}\circled{3}}
    {\circled{2}\circled{4}}
\end{problem}

\begin{problem}
如图，在圆心角为直角的扇形 OAB 中，分别以 OA, OB 为直径作两个半圆. 在扇形 OAB 内随机取一点，则此点取自阴影部分的概率是（\quad）

\choices
    {\(1 - \frac{2}{\pi}\)}
    {\(\frac{1}{2} -\frac{1}{\pi}\)}
    {\(\frac{2}{\pi}\)}
    {\(\frac{1}{\pi}\)}
\end{problem}

\begin{problem}
函数 \( f(x) = x\cos x^2 \) 在区间 \([0, 4]\) 上的零点个数为（\quad）

\choices
    {4}
    {5}
    {6}
    {7}
\end{problem}

\begin{problem}
我国古代数学名著《九章算术》中“开立圆术”曰：置积尺数，以十六乘之，九而一，所得开立方器，正切圆径. “开立圆术”相当于给出了已知球的体积 \(V\) ，求其直径 \(d\) 的一个近似公式 \(d \approx \sqrt[3]{\frac{16}{3} V}\) . 人们运用过一些类似的近似公式. 根据 \(\pi = 3.14159 \cdots\) 判断，下列说法众尺中最精确的是（\quad）

\choices
    {\(d \approx \sqrt[3]{\frac{16}{9} V}\)}
    {\(d \approx \sqrt[3]{2V}\)}
    {\(d \approx \sqrt[3]{\frac{200}{151}}\)}
    {\(d \approx \sqrt[3]{\frac{211}{11}}\)}
\end{problem}

\section{填空题}

\begin{problem}
设 \(\triangle ABC\) 的内角 \(A, B, C\) 所对的边分别为 \(a, b, c\). 若 \((a + b - c)(a + b + c) = ab\)，则 \(C = \fillinblank{}\).
\end{problem}

\begin{problem}
阅读如图所示的程序框图，运行相应的程序，输出的结果是 \fillinblank{} \fillinblank{}.
\end{problem}

\begin{problem}
回文数是指从左到右读与从右到左读都一样的正整数. 如 22, 121, 3443, 94249 等. 显然 2 位回文数有 9 个：11, 22, 33，…，99, 3 位回文数有 90 个：101, 111, 121, …, 191, 202，…，999. 则 (1) 4 位回文数有 个; (2) \( 2n+1 \) ( \( n \in N^{*} \) ) 位回文数有 个 \fillinblank{}.
\end{problem}

\begin{problem}
如图，双曲线 \(\frac{x^2}{a^2} - \frac{y^2}{b^2} = 1 (a, b > 0)\) 的两顶点为 \(A_1, A_2\)，虚轴两端点为 \(B_1, B_2\)，两焦点为 \(F_1, F_2\). 若以 \(A_1A_2\) 为直径的圆内接于菱形 \(F_1B_1B_2B_2\)，切点分别为 \(A, B, C, D\). 则 (1) 双曲线的离心率 \(e = \fillinblank{}\). (1) 双曲线的离心率 \( e = \) ； (2) 菱形 \(F_{1}B_{1}F_{2}B_{2}\) 的面积 \(S_{1}\) 与矩形ABCD的面积 \(S_{2}\) 的比值 \(\frac{S_{1}}{S_{2}} = \fillinblank{}\).
\end{problem}

\begin{problem}
如图，点 \(D\) 在 \(\odot O\) 的弦 \(AB\) 上移动，\(AB = 4\); 连接 \(OD\)，过点 \(D\) 作 \(OD\) 的垂线交 \(\odot O\) 于点 \(C\)，则 \(CD\) 的最大值为 \fillinblank{} \fillinblank{}.
\end{problem}

\begin{problem}
在直角坐标系 \(xOy\) 中，以原点 \(O\) 为极点，\(x\) 轴的正半轴为极轴建立极坐标系. 已知格线 \(\theta = \frac{\pi}{4}\) 与曲线 \(\left\{ \begin{array}{l} x = t + 1, \\ y = (t - 1)^2, \end{array} \right.\) (\(t\) 为参数) 相交于 \(A, B\) 两点，则线段 \(AB\) 的中点的直角坐标为 \fillinblank{} \fillinblank{}.
\end{problem}

\section{解答题}

\begin{problem}
已知向量 \( \bs{a} = (\cos \omega x - \sin \omega t, \sin \omega t) \) ，\( \bs{b} = (-\cos \omega t - \sin \omega t, 2\sqrt{3} \cos \omega t) \) ，设函数 \( f(x) = \bs{a} \cdot \bs{b} + \lambda (x \in \R) \) 的图象关于直线 x = 7 对称，其中 \( \omega, \lambda \) 为常数，且 \( \omega \in \left(\frac{3}{2}, 1\right) \) . (1) 求函数 \( f(x) \) 的最小正周期; (2) 若 \( y = f(x) \) 的图象经过点 \( \left(\frac{\pi}{4}, 0\right) \)，求函数 \( f(x) \) 在区间 \( \left[0, \frac{3\pi}{5}\right] \) 上的取值范围.
\end{problem}

\begin{problem}
已知等差数列 \(\{n_{n}\}\) 前三项的和为\(-3\)，前三项的积为8. (1) 求等差数列 \( \{n_{n}\} \) 的通项公式; (2) 若 \(a_{2}, a_{3}, a_{1}\) 成等比数列，求数列 \(\{\{n_{n}\}\}\) 的前 \(n\) 项和.
\end{problem}

\begin{problem}
根据以往的经验，某工程施工期间的降水量 X (单位: mm) 对工期的影响如下表: 降水量X X\&lt;300 300≤X\&lt;700 700≤X\&lt;900 X≥900 工期延误天数Y 0 2 6 10 历年气象资料表明，该工程施工期间降水量 X 小于 300, 700, 900 的概率分别为 0.3, 0.7, 0.9. 求: (1) 工期延误天数 Y 的均值与方差; (2) 在降水量 X 至少是 300 的条件下，工期延误不超过 6 天的概率.
\end{problem}

\begin{problem}
(1) 已知函数 \( f(x) = rx - x^2 + (1 - r)(x > 0) \)，其中 \( r \) 为有理数，且 \( 0 < r < 1 \)，求 \( f(x) \) 的最小值； (2) 试用 (1) 的结果证明如下命题: 设 \(a_{1} \geqslant 0, a_{2} \geqslant 0, b_{1}, b_{2}\) 为正有理数. 若 \(b_{1} + b_{2} = 1\)，则 \(a_{1}^{b_{1}} a_{2}^{b_{2}} \leqslant a_{1} b_{1} + a_{2} b_{2}\); (3) 请将 (2) 中的命题推广到一般形式，并用数学归纳法证明你所推广的命题. 注：当 \(\alpha\) 为正有理数时，有求导公式 \((x^{\alpha})^{\prime} = \alpha x^{\alpha -1}\)
\end{problem}

\begin{problem}
如图 1, \(\angle ACB = 45^{\circ}\)，\(BC = 3\)，过动点 \(A\) 作 \(AD \perp BC\)，垂足 \(D\) 在线段 \(BC\) 上且异于点 \(B\)，连接 \(AB\). 沿 \(AD\) 将 \(\triangle ABD\) 折起，使 \(\angle BDC = 90^{\circ}\) (如图 2 所示). (1) 当 \(BD\) 的长为多少时，三棱锥 \(A - BCD\) 的体积最大; (2) 当三棱锥 \(A - BCD\) 的体积最大时，设点 \(E, M\) 分别为棱 \(BC, AC\) 的中点，试在棱 \(CD\) 上确定一点 \(N\)，使得 \(EN \perp BM\)，并求 \(EN\) 与平面 \(BMN\) 所成角的大小.

图1

图2
\end{problem}

\begin{problem}
设 \(A\) 是单位圆 \(x^{2} + y^{2} = 1\) 上的任意一点，\(l\) 是过点 \(A\) 与 \(x\) 轴垂直的直线，\(D\) 是直线 \(|l\) 与 \(x\) 轴的交点，点 \(M\) 在直线 \(l\) 上，且满足 \(\|DM\| = m\|DA\|\) (\(m > 0\) 且 \(m \neq 1\)). 当点 \(A\) 在圆上运动时，记点 \(M\) 的轨迹为曲线 \(C\). (1) 求曲线 \(C\) 的方程，判断曲线 \(C\) 为何种圆锥曲线. 并求其焦点坐标; (2) 过原点且斜率为 \(k\) 的直线交曲线 \(C\) 于 \(P, Q\) 两点，其中 \(P\) 在第一象限，它在 \(y\) 轴上的射影为点 \(N\)，直线 \(QN\) 交曲线 \(C\) 于另一点 \(H\). 是否存在 \(m\)，使得对任意的 \(k > 0\)，都有 \(PQ \perp PH^2\) 若存在，求 \(m\) 的值; 若不存在，请说明理由.
\end{problem}

